\documentclass[11pt,a4paper]{article}

\usepackage[T1]{fontenc}
\usepackage[utf8]{inputenc}
\usepackage[a4paper,margin=1.8cm]{geometry}
\usepackage{amsmath,amssymb}
\usepackage{booktabs,longtable,array,tabularx}
\usepackage{graphicx}
\usepackage{microtype}
\usepackage{xcolor}
\usepackage{xr-hyper}
\usepackage{hyperref}
\externaldocument[M-]{main}

\hypersetup{colorlinks=true,linkcolor=blue,citecolor=blue,urlcolor=blue}

\graphicspath{{./figures/}}

\renewcommand{\thesection}{S\arabic{section}}
\renewcommand{\thesubsection}{S\arabic{section}.\arabic{subsection}}
\renewcommand{\thetable}{S\arabic{table}}
\renewcommand{\thefigure}{S\arabic{figure}}
\newcommand{\code}[1]{\texttt{#1}}
\newcommand{\modelname}{\texttt{minimocas}}
\newcommand{\todo}[1]{\textcolor{red}{\textbf{[TODO: #1]}}}
\setlength{\tabcolsep}{4pt}
\renewcommand{\arraystretch}{1.12}

\title{Supplementary Material for\\
\emph{Large-scale traffic flow reconstruction from heterogeneous mobility sensors}}
\author{Federico Bellisardi, Chiara Mizzi, Alessandro Fabbri and Armando Bazzani}
\date{}

\begin{document}
\maketitle


\noindent This Supplementary Material documents the data layers, the simulation model, the route-choice variants, the detector correction and additional validation results of the main manuscript, and lists the commands that reproduce them.

\section{Data layers}
\label{sec:data-layers}

\begin{table}[htbp]\centering
\caption{Data layers and their role in the reconstruction.}
\label{tab:layers}
\begin{tabularx}{\textwidth}{@{}lXXl@{}}\toprule
Layer & Content & Role & Period \\\midrule
Mobile-phone OD (TIM) & hourly device movements between census areas, aggregated to municipalities and provinces & demand & January 2020 \\
Road graph & 3,925 directed polylines, 2,548 nodes, turn restrictions, derived from OpenStreetMap & network & -- \\
OSM road classes & motorway to unclassified, nearest ways within 30~m & route-choice parameters & extract of 26/10/2026 \\
MTS loop detectors & 15-min counts per station and direction & local correction; volume validation & Feb 2020, Jun--Jul 2021 \\
Viasat GPS, link counts & hourly probes per matched link & profile validation & 10--17/02/2020, 28/06--04/07/2021 \\
Viasat GPS, raw traces & private-vehicle fixes in Forl\`i--Cesena & speeds and class usage for route choice & same periods \\
\bottomrule\end{tabularx}
\end{table}

The OD tables are hourly. Their labels refer to the start of the hour in local time: the weekday OD profile of 10--14 February 2020 and the mean profile of the 90 detectors complete on those days have correlation 0.96 at zero lag, 0.94 with the OD profile shifted one hour earlier and 0.74 with it shifted one hour later. Demand is generated in four classes: observed source and destination (\code{src-dest}, municipalities, and \code{scr\_dest\_prov}, provinces), observed source and random destination (\code{src-rnd}), random source and observed destination (\code{rnd-dest}) and random source and destination (\code{rnd-rnd}). On a weekday, the random-endpoint classes account for 66\% of the generated agents (\code{rnd-rnd} 33\%, \code{src-rnd} 22\%, \code{rnd-dest} 12\%).

MTS counts of 2021 were obtained from the monthly open-data files of Regione Emilia-Romagna (dataset DS-1944, CC BY), summing all vehicle classes. The open-data time labels are in local time; converted to UTC they reproduce the counts of the February 2020 archive used for 2020 (ratio of totals 0.995 on 12 February 2020).

\section{Simulation model}
\label{sec:model}

Agents move with the speed of their link, $v_i$, reduced linearly with density $\rho_i$ below the critical density and set to a critical speed above it. Link speeds are read in km/h from the cartography and assigned by road class from GPS (main text, Table~\ref{M-tab:classes}). Random-endpoint agents have a trip-time budget of 5,600~s; source--destination agents follow the best-path matrix of their type. When the turn-restricted best-path search finds no feasible next node for the cost of their type, the agent uses the shortest-distance path matrix; this happened for 50 next-node decisions in two of the fifteen simulated days. The simulator is deterministic (fixed random seed). A simulated day takes 9--12 minutes on an 8-core workstation running four days in parallel.

\section{Route-choice variants}
\label{sec:routing}

Table~\ref{tab:routing} lists the variants simulated for 10 February 2020. Speeds alone (V1--V3) do not change the distribution of traffic over road classes, because two thirds of the agents follow transition probabilities rather than best paths. Transition weights by road class (V4--V6) raise the share of main roads and the agreement with GPS counts in the calibration area. The selected variant combines the usage cost of Eq.~\eqref{M-eq:route_cost} with class transition weights. Correlations with the detectors of a single day (53 coils) are noisy and were not used for selection; the comparison over full campaigns is in the main text. Table~\ref{tab:routing_shape} shows that route choice hardly changes the hourly profiles.

\input{tables/routing_variants}
\input{tables/routing_shape}

\section{Detector correction}
\label{sec:correction}

The correction starts from the residual of Eq.~\eqref{M-eq:coil_residual} on each coil link and hour. The ratio $r_i(h)$ converts the coil count, a count of passages at a section, into the scale of the simulated unique-link count, which counts distinct agents present on the whole polyline during the hour. Downstream, the increment of a link is passed to each admissible outgoing link multiplied by that link's share of the simulated outgoing counts; upstream, the same rule is applied to incoming links. Each branch carries its own increment and its own path length. Propagation stops when the increment is below one vehicle, after three polylines along the path, at another coil link, at a forbidden turn or at a closed polyline. Every directed polyline of the raw simulation is kept in the output, and the hours of the observed table are aligned with the local-time hours of the simulation.

Table~\ref{tab:loo} summarizes the leave-one-detector-out test, Table~\ref{tab:depth} the dependence on the propagation depth and Table~\ref{tab:distance} the effect of the correction on the hourly profiles of GPS links by network distance from the nearest coil.

\input{tables/loo}
\input{tables/depth}
\input{tables/distance}

\section{Additional validation results}
\label{sec:validation}

Table~\ref{tab:probes} reports the profile error by number of probes per link. The error decreases with the number of probes for both the reconstruction and the MTS profile, as expected when part of the difference is sampling noise of the probe counts. Matched links with zero simulated or GPS total are excluded; absent link-hour records of the GPS extract are counted as zero.

\input{tables/probe_strata}

\section{Reproducibility}
\label{sec:reproducibility}

All results are produced by the scripts in \code{article1/analysis} of the code repository, run from its root:
\begin{enumerate}
\item \code{od\_clock\_check.py}: hour convention of the OD tables;
\item \code{carto\_geojson.py}, \code{osm\_road\_class.py}, \code{gps\_link\_speeds.py}: geometry, road classes and GPS link speeds;
\item \code{routing\_variants.py prepare|run|evaluate}: route-choice variants and selection;
\item \code{run\_corrected\_sims.py}: simulations of the validation days;
\item \code{build\_coils\_table.py}: 2021 detector table from the open-data files;
\item \code{reconstruct\_and\_validate.py assimilate|validate|loo}: correction, GPS validation, leave-one-out;
\item \code{volume\_check.py}, \code{compare\_routing.py}, \code{distance\_to\_detector.py}, \code{sensitivity\_fixed.py}: volume tests, routing comparison, reach of the correction, depth sensitivity;
\item \code{make\_figures.py}, \code{make\_supp\_tables.py}: figures and tables.
\end{enumerate}
Proprietary inputs (mobile-phone OD tables, Viasat link counts and traces) are required for steps 1--6 and are not distributed.

\end{document}
